\chapter{Conclusion and Future Work}
\label{ch:conclusion}

\section{Conclusion}

This thesis presented a systematic, simulation-based evaluation of two prominent routing
protocols for Wireless Mesh Networks --- the reactive AODV (RFC~3561) and the proactive
DSDV (Perkins \& Bhagwat~1994) --- using OMNeT++~6.3 and the INET~4.5 framework.
Simulations were conducted across five network densities ($N \in \{10, 20, 30, 40, 50\}$
nodes) in a $1000\,\text{m} \times 1000\,\text{m}$ area with IEEE~802.11g radio links
and a many-to-one UDP traffic pattern representative of IoT data collection.

The key conclusions are:

\begin{enumerate}
  \item \textbf{Performance crossover at approximately $N = 20$.} DSDV offers lower
    end-to-end delay at $N = 10$ due to pre-populated routing tables, but AODV
    outperforms DSDV on all five metrics (throughput, packet count, delay, energy, and
    drop rate) at $N \geq 30$. The crossover is driven by a shift in the dominant
    performance bottleneck from route availability (sparse) to channel capacity (dense).

  \item \textbf{AODV delivers 35\% higher throughput at $N = 50$.} AODV achieves
    $\approx 2.5$~MB of delivered data compared with $\approx 1.85$~MB for DSDV. The
    difference is caused by DSDV's periodic routing broadcasts consuming channel airtime
    that AODV reserves for application data.

  \item \textbf{AODV achieves zero routing failures across all densities.} DSDV suffers
    24--31 ``no route found'' drops at $N = 10$ (convergence lag) and 121+ queue-overflow
    drops at $N = 50$ (MAC saturation). For reliability-critical IoT applications,
    AODV's zero-drop record is a decisive advantage.

  \item \textbf{MAC-layer evidence confirms congestion collapse at $N = 50$ for DSDV.}
    Contention-window traces show CW escalating to 511--1023 slots, with repeated
    7-retry aborts. This is direct physical evidence that DSDV's performance degradation
    is not an algorithmic failure but a channel saturation effect triggered by the MAC
    competition between routing broadcasts and data traffic.

  \item \textbf{AODV is 8\% more energy-efficient at $N = 50$.} The energy advantage
    arises from AODV's on-demand control traffic model, which keeps the radio in the
    transmit and receive states less frequently than DSDV's periodic broadcasts. This
    advantage is likely to be larger in longer operational scenarios.
\end{enumerate}

\section{Deployment Guidelines}

Based on the simulation results, the following protocol-selection guidelines are
recommended for WMN practitioners:

\subsection{When to Choose DSDV}

DSDV should be preferred when:

\begin{itemize}
  \item \textbf{Network size is small ($N < 20$ nodes).} Below the crossover density,
    DSDV's pre-populated routes deliver substantially lower latency.
  \item \textbf{Latency is the primary concern} and reliability can be sacrificed
    occasionally. Home automation systems, hospital patient monitoring (where a missed
    reading can be tolerated), and similar small-scale deployments fall into this category.
  \item \textbf{Communication is frequent.} If every pair of nodes communicates
    constantly, AODV's route-discovery overhead is amortised less effectively, and
    DSDV's always-ready routes become an advantage.
\end{itemize}

\subsection{When to Choose AODV}

AODV should be preferred when:

\begin{itemize}
  \item \textbf{Network size is large ($N > 20$ nodes).} Above the crossover, AODV's
    on-demand control model prevents channel saturation.
  \item \textbf{Reliability is critical.} Industrial IoT, smart grid metering, and
    emergency communications cannot tolerate sustained packet-drop episodes. AODV's
    zero-drop record makes it the safer choice.
  \item \textbf{Energy efficiency matters.} Battery-powered sensor nodes on a WMN with
    $N > 20$ will achieve longer operational lifetimes with AODV.
  \item \textbf{Many-to-one traffic is the dominant pattern.} IoT data collection
    scenarios create hot-spots near the sink. AODV's reactive mechanism naturally
    concentrates routes near active paths, avoiding the broadcast overhead that saturates
    the hot-spot region under DSDV.
\end{itemize}

\subsection{Decision Matrix}

Table~\ref{tab:decision_matrix} provides a concise decision matrix for protocol selection.

\begin{table}[ht]
\centering
\caption{Protocol selection decision matrix for WMN deployments.}
\label{tab:decision_matrix}
\begin{tabular}{@{}lll@{}}
\toprule
\textbf{Criterion} & \textbf{Prefer AODV} & \textbf{Prefer DSDV} \\
\midrule
Network size         & $N > 20$                      & $N < 20$ \\
Primary requirement  & Throughput / reliability      & Low latency \\
Traffic pattern      & Many-to-one, infrequent       & All-pairs, frequent \\
Energy budget        & Constrained (battery)         & Unconstrained (mains) \\
MAC protocol         & 802.11 (any)                  & Low-bitrate IoT radios \\
Mobility             & Mobile nodes                  & Stationary / low-mobility \\
Fault tolerance      & Critical (zero drops required)& Tolerant (occasional drops) \\
Deployment example   & Smart city, industrial IoT    & Home automation, hospital \\
\bottomrule
\end{tabular}
\end{table}

\section{Future Work}

The following extensions would strengthen and broaden the conclusions of this thesis:

\paragraph{Node mobility.}
This study used stationary nodes. Introducing mobility (random waypoint, Gauss-Markov)
would activate AODV's RERR/re-discovery mechanism and DSDV's rapid-update response,
potentially shifting the crossover density.

\paragraph{Additional protocols.}
Including OLSR~\cite{clausen2003olsr} (which uses multipoint relays to reduce broadcast
flooding overhead compared with DSDV) and HWMP (the protocol mandated by IEEE~802.11s
for standardised mesh networking) would provide a more comprehensive comparison.

\paragraph{Heterogeneous traffic.}
The current study uses a single CBR source pattern. Real IoT deployments mix periodic
sensor readings with event-driven alarm messages and multimedia streams. Evaluating AODV
and DSDV under mixed traffic classes, including TCP flows with congestion control, would
provide more practically relevant results.

\paragraph{Energy-aware routing.}
Routing decisions could incorporate residual energy as a metric, preferring paths through
nodes with higher remaining battery capacity. This would extend network lifetime and
distribute the energy burden more evenly. Such a mechanism is straightforward to
implement in INET by modifying the routing metric computation.

\paragraph{Variable MAC parameters.}
The retry limit, $\text{CW}_{\min}$, and $\text{CW}_{\max}$ significantly affect the
crossover density. A parametric study varying these values would characterise the
sensitivity of the crossover to MAC configuration.

\paragraph{Physical hardware validation.}
Simulation results should be validated on a physical testbed (e.g., Raspberry Pi~4
nodes running OpenWRT with batman-adv or babeld) to account for real-world effects
including RF interference, multi-path fading, asymmetric link quality, and CPU
scheduling latency.

\paragraph{Security analysis.}
WMN routing protocols are vulnerable to attacks including black-hole attacks (a malicious
node advertises routes to attract traffic and then drops it), wormhole attacks, and Sybil
attacks. Extending the evaluation to include adversarial node behaviour would add
critical practical relevance, particularly for the IoT security domain.

\section{Final Remarks}

This project demonstrated that simulation is an effective and reproducible method for
evaluating the trade-offs between routing protocols in Wireless Mesh Networks. The
OMNeT++/INET platform provided a high-fidelity model of the IEEE~802.11 MAC and
physical layers, yielding results that closely reflect the mechanisms observed in
real-world 802.11 networks.

The central finding --- that the dominant performance bottleneck shifts from route
availability to channel capacity as network density increases --- is not specific to
AODV and DSDV. It is a general principle that applies to any comparison between reactive
and proactive routing protocols in any wireless network that uses a shared-medium MAC
protocol. This principle should guide protocol designers and network engineers alike:
the best routing protocol is the one whose control overhead is best matched to the
channel capacity available at the target deployment density.

Wireless Mesh Networks continue to play a growing role in smart city infrastructure,
industrial automation, agricultural monitoring, and emergency communications. As
deployments grow larger and energy constraints become more stringent, the quantitative
understanding provided by studies such as this one becomes increasingly important for
making informed protocol-selection decisions.
